% Formula and ink-to-index acceptance fixture for issue 4555.
% K1: a dashed enclosure states that its member tensor bodies form one
% blocked local object.  Virtual and physical legs remain external indices;
% the enclosure is presentation ink, never a contraction.
% The five planned consumers reduce to: one long group, several disjoint
% pairs, a two-site capability, and two three-site blocked rows.  The cited
% source for TNMPDOTwoSiteBlocking draws no enclosure, so this file exercises
% the general two-site capability without assigning that extra ink to it.
% K7: a region encloses the measured ink of named free atoms and named joins.
% Its boundary denotes a selected/grouped subgraph and introduces no index.
\documentclass[border=6pt]{standalone}
\usepackage{amsmath}
\usepackage{tenkz}

\begin{document}
\begin{minipage}{16.5cm}
\centering

% A^{[L]}_{i_1\ldots i_L}=A^{i_1}\cdots A^{i_L}; the box groups
% the four displayed tensor bodies while every open leg stays an index.
\[
\begin{tenkz}[physical=up, tensor style=box]
  \tn[up=$i_1$]{A}\tnspan[box, label pos=west]{4}{A^{[L]}} &
  \tn[up=$i_2$]{A} & \tndots & \tn[up=$i_L$]{A}
\end{tenkz}
\]

% Two disjoint length-two blocked factors; neither box captures its neighbour.
\[
\begin{tenkz}[tensor style=box]
  \tn{K}\tnspan[box, label pos=south]{2}{B_1} & \tn{K} &
  \tn{K}\tnspan[box, label pos=south]{2}{B_2} & \tn{K}
\end{tenkz}
\qquad
% Abstract two-site capability only: no claim about MPDO_K=MM.png.
\begin{tenkz}[physical=updown, tensor style=box]
  \tn{K}\tnspan[box, label pos=north]{2}{K^{[2]}} & \tn{K}
\end{tenkz}
\]

% Asymmetric row lengths test independent three-site blocked factors.
\[
\begin{tenkz}[rows={op,op}, tensor style=box]
  \tn{A_1}\tnspan[box, label pos=west]{3}{A'} & \tn{A_2} & \tn{A_3} \\
  \tn{C_1}\tnspan[box, label pos=east]{2}{C'} & \tn{C_2} &
\end{tenkz}
\]

% The named joins j_0,j_1 carry summed virtual indices.  R is their selected
% support together with U_0,U_1; P groups the rank-one bond object phi.
\[
\begin{tenkzfree}
  \tnput[box, ports={south:virtual}]{u0}{(0,1.1)}{U_0}
  \tnput[box, ports={south:virtual}]{u1}{(2.5,1.4)}{U_1}
  \tnput[dot, ports={west:virtual,east:virtual},
         label pos=south]{phi}{(1.2,-0.2)}{\varphi}
  \tnjoin[name=j0, route=vh]{u0.south}{phi.west}
  \tnjoin[name=j1, route=hv]{u1.south}{phi.east}
  \tnregion[group, label={$P$}, label pos=south]{phi}
  \tnregion[slot=selected, outline, label={$R$}, label pos=north,
            name=R]{u0,u1,j0,j1}
\end{tenkzfree}
\qquad
% Oblique, asymmetric atom labels are part of each measured member box.
\begin{tenkzfree}
  \tnput[dot, label pos=west]{a}{(0,0)}{A_{\alpha}}
  \tnput[pill]{b}{(1.6,0.8)}{B^{\beta}}
  \tnregion[slot=secondary, label={$S$}, label pos=north east]{a,b}
\end{tenkzfree}
\]

% No-feature control: loading the enclosure implementation does not add ink to
% an ordinary grid.  The regression test compares two copies structurally.
\[
\begin{tenkz}[tensor style=box]
  \tn{X} & \tn{Y}
\end{tenkz}
\]

\end{minipage}
\end{document}
